\section{Related Works}\label{planta:sec:related}

\subsection{Self-Adaptive Robotic Systems}\label{planta:sec:related_sas}

The mapping study on ``robotics software architecture-based self-adaptive systems''~(RSASSs) by Alberts et al.~\cite{ALBERTS2025112258} identified 37 primary studies published since 2011. Among these, only four works~\cite{park2012task, runtime_variability_lotz, rra_nico, CAMARASoftwareArchitectureTask2020} explicitly address, to varying extents, the dependencies between task execution and architectural adaptation. Additionally, the two studies~\cite{braberman2018extended, temoto} also explore this relationship,  and the ROSA framework proposes a systematic solution to TACA~\cite{rosa}. However, only the approaches~\cite{braberman2018extended, CAMARASoftwareArchitectureTask2020,rosa} fully address TACA\@.
With respect to knowledge-based methods, Alberts et al.\ identified six works~\cite{park2012task, stefan_adaptive_runtime_models, graph_based_nico, p20_stefan_ice, p26_darko_jasper, suave} that utilize knowledge representation techniques to encode the information required for adaptation. From these, the works~\cite{p26_darko_jasper, suave} focus on demonstrating the use of Metacontrol~\cite{hernandez2018self, bozhinoski2022mros} in specific robotic applications rather than proposing complete solutions for RSASSs.
Concerning planning, using PDDL-based planners for RSASSs is still mostly unexplored. Alberts et al.\ identified only the approach \enquote{Model Based Diagnosis (MBD)}~\cite{zaman2013integrated, zaman2019fault}, which employs PDDL-based planning to generate architectural adaptation plans, without explicitly targeting TACA\@.

Cámara et al.~\cite{CAMARASoftwareArchitectureTask2020} proposed a solution to TACA that generates an optimal policy for task execution and architectural configuration. Their method is decomposed into four steps: (i) generate all legal architectural configurations based on the robot’s architecture model, (ii) generate all possible reconfiguration plans between configurations, (iii) generate the \emph{N} best task plans to accomplish the robot's mission goal, (iv) compute an optimal policy for task execution and configuration, given the reconfiguration and task plans. Their experiments show that the generation time is impractical: Generating an optimal task and configuration policy for an AGV operating in a graph-like environment with 58 nodes and with six architectural configurations available takes, on average, $15.1$ seconds.

Silva et al.~\cite{rosa} proposed ROSA as a solution to TACA, focusing on providing a reference framework for robot self-adaptation that promotes reusability, composability, and extensibility. A limitation of ROSA is that it only supports reactive adaptation of the robot's configuration. ROSA determines the robot's configuration for each action only when it starts executing, and it does not produce a complete configuration plan ahead of time.

To address the limitations of these approaches, this work proposes an ontology-mediated planning approach that enables the timely computation of a satisficing task and configuration plan without relying on probabilistic models of real-world uncertainties.
% In conclusion, the method proposed by Cámara et al.~\cite{CAMARASoftwareArchitectureTask2020} is, to the best of our knowledge, the only approach that plans task actions and reconfiguration actions jointly and produces a complete plan ahead of execution. However, it presents practical limitations: it relies on prior knowledge of the probabilities associated with real-world uncertainties, which are often difficult to obtain in practice. Moreover, due to its focus on computing optimal solutions, the method is not fast enough for reactive use in many real-world robotics scenarios, where timely decision-making is critical. To address these limitations, this work proposes an ontology-mediated planning approach that enables the timely computation of a satisficing task and configuration plan without relying on probabilistic models of real-world uncertainties.


\subsection{Combining Ontologies and Planning}\label{planta:sec:related_ontology_planning}

\paragraph{Knowledge and Action Bases.}
Several works have explored combining ontologies with planning problems. The approach proposed in this work is most closely related to the work on \emph{Knowledge and Action Bases (KABs)}~\cite{DBLP:journals/jair/HaririCMGMF13}, where an ontology is used to provide additional insights about a planning state by using queries to the ontology in the planning domain. The actions in the planning domain describe changes to the ontology over time.

There are different approaches to generate a plan for KABs, namely (i) rewriting ontology queries to complex preconditions~\cite{DBLP:conf/ijcai/CalvaneseMPS16}, (ii) rewriting ontology axioms to derived predicates~\cite{DBLP:conf/aaai/Borgwardt0KKNS22}, and (iii) enriching the planning domain with grounded answers to the ontology queries~\cite{DBLP:conf/ecai/JohnK24}. All the approaches have in common that they perform planning as a two-step process: First, the ontology is analyzed and a new planning domain is created that encodes the axioms from the ontology. Secondly, an off-the-shelf planner is used to find a plan using the new domain and the original problem. The first approach only works for axioms in DL-Lite~\cite{DBLP:conf/ijcai/CalvaneseMPS16} and leads to rather long planning times~\cite{DBLP:conf/aaai/Borgwardt0KKNS22}, despite optimizations in the expansion algorithm~\cite{DBLP:conf/kr/Borgwardt0KS21}. The second approach requires a Horn-Logic fragment of the ontology axioms, but it is, in general, the fastest among the three~\cite{DBLP:conf/aaai/Borgwardt0KKNS22, DBLP:conf/ecai/JohnK24}. The third approach supports all kinds of axioms, but it is slower than the second approach, especially in instances where finding a plan is relatively simple~\cite{DBLP:conf/ecai/JohnK24}.

The approach presented in this work is most similar to the second approach, where ontology axioms are encoded as derived predicates in the planning formulation. However, due to the special requirements of TACA, a custom translation from axioms to derived predicates is used.

\paragraph{Ontologies as Modeling Languages.} Some studies have explored the use of ontologies as a modeling language for planning problems in the robotics domain~\cite{DBLP:journals/ras/BalakirskyKKPSG13,DBLP:journals/ewc/BorgoCOU19}. These efforts try to make modeling of complex domains easier by providing frameworks for complex classes and their relations. It is a known limitation of PDDL that modeling domains with complex types directly in PDDL is not well supported~\cite{GOLDMANpddlTypeProblems2012}. This work adopts a related approach by encoding PDDL type information within the ontology to support semantically rich planning representations.


%%% Local Variables:
%%% mode: latex
%%% TeX-master: "main"
%%% End:
